% Abhakseite „Das kann ich“ – Zone nur im Gesamt (\mitzone)
%% TODO Ich-kann-Sätze: Platzhalter wie in den Aufgabendateien
\begin{abhakseite}
\ifdefined\mitzone
\abhakgruppe{Kennst du schon}
\abhak{1}{Addieren und Subtrahieren negativer Zahlen}
\abhak{2}{Multiplizieren mit Vorzeichen}
\abhak{3}{Dezimalzahlen und einfache Brüche als Vorzahlen}
\abhak{4}{Punkt vor Strich mit Zahlen}
\abhak{5}{Addieren und Subtrahieren negativer Zahlen – Fehler finden}
\abhak{6}{Addieren und Subtrahieren negativer Zahlen – selbst rechnen}
\fi
\abhakgruppe{Einheit 1 · Terme aufstellen und berechnen}
\abhak{7}{Term aufstellen}
\abhak{8}{Termwert berechnen}
\abhak{9}{Situation zu Term angeben}
\abhak{10}{Term aufstellen – Fehler finden, Darstellungswechsel, Anwendung}
\abhakgruppe{Einheit 2 · Terme zusammenfassen}
\abhak{11}{Gleichartige Glieder erkennen}
\abhak{12}{Vorzahl lesen}
\abhak{13}{Vorzeichen als Teil des Gliedes}
\abhak{14}{Zusammenfassen}
\abhak{15}{Malnehmen}
\abhak{16}{Zusammenfassen – Fehler finden, Begründen, Darstellungswechsel, Anwendung}
\abhakgruppe{Einheit 3 · Klammern auflösen}
\abhak{17}{Zeichen vor der Klammer feststellen}
\abhak{18}{Klammern}
\abhak{19}{Klammern – Fehler finden, Begründen, Darstellungswechsel}
\abhakgruppe{Einheit 4 · Ausklammern}
\abhak{20}{Ausklammern}
\abhak{21}{Ausklammern – Fehler finden}
\end{abhakseite}
